% GENERATED FILE. Edit canonical lessons, then run npm run generate:books.
% canonical-source-hash: fnv1a64:8f5b7b3f0ee578b3
% canonical-lessons: MW-C118-chulho, MW-C118-tavo, MW-C118-karahi, MW-C118-hadi, MW-C118-gharo

\chapter{Hearth, Griddle, Wok, Cooking pot, Water pot}
\label{ch:mw-hearth-griddle-wok-cooking-pot-water-pot}
\hlchaptermodality{118}

\begin{chapteropening}
\textbf{By the end of this chapter:} \emph{I can say a hearth, a griddle, a wok, a cooking pot and a water pot.}

Complete the last lesson of chapter 118: I can say a hearth, a griddle, a wok, a cooking pot and a water pot.
\end{chapteropening}

\section[\texorpdfstring{\mw{चूल्हो} (chūlho)}{चूल्हो (chūlho)}]{\mw{चूल्हो} (chūlho) — a hearth, a stove}

\label{lesson:MW-C118-chulho}



\begin{quote}
Before the new one: say the Marwadi for because of, then the Marwadi for because.
\end{quote}

\subsection*{You'll want to know: \mw{चूल्हो}}
\textbf{\mw{चूल्हो}} — \emph{chūlho} — \textquotedblleft{}a hearth, a stove\textquotedblright{}.

\begin{cousinweb}
\textbf{\mw{चूल्हो}} (\emph{chūlho}) is a tadbhava, a word inherited from Sanskrit and reshaped by sound change. It comes from \emph{cullī}, \textquotedblleft{}fireplace\textquotedblright{}. CDIAL gives no Marwari form in this entry (CDIAL 4879); the match is through Hindi \emph{cūlhā}. CDIAL takes the Sanskrit word itself from Dravidian.
\end{cousinweb}

\begin{grammarlens}[title={What you've built}]
One more word for this chapter.
\end{grammarlens}

\subsection*{Your turn}
\begin{itemize}
\raggedright
  \item \emph{Say it:} \emph{chūlho}
  \item \emph{Say it:} \emph{chūlho}, once more
  \item \emph{Say it:} say \emph{chūlho}
  \item \emph{Recall:} say the Marwadi for because of, then the Marwadi for because, then say \emph{chūlho} again
\end{itemize}

\begin{tcolorbox}[breakable,colback=teal!4,colframe=teal!35!black,title={Before you move on}]
What does \mw{चूल्हो} mean? (\textquotedblleft{}A hearth, a stove\textquotedblright{}.) Say it once more.
\end{tcolorbox}
\section[\texorpdfstring{\mw{तवो} (tavo)}{तवो (tavo)}]{\mw{तवो} (tavo) — a griddle}

\label{lesson:MW-C118-tavo}



\begin{quote}
Before the new one: say the Marwadi for because, then the Marwadi for a hearth.
\end{quote}

\subsection*{You'll want to know: \mw{तवो}}
\textbf{\mw{तवो}} — \emph{tavo} — \textquotedblleft{}a griddle\textquotedblright{}.

\begin{cousinweb}
\textbf{\mw{तवो}} (\emph{tavo}) is a tadbhava, a word inherited and reshaped by sound change. It comes from \emph{*tapaka}, \textquotedblleft{}frying pan\textquotedblright{}, a form CDIAL reconstructs from its descendants. CDIAL gives no Marwari form in this entry (CDIAL 5670); the match is through Hindi \emph{tawā} and Gujarati \emph{tavo}. It is built on the root \emph{tap}, \textquotedblleft{}to heat\textquotedblright{}.
\end{cousinweb}

\begin{grammarlens}[title={What you've built}]
One more word for this chapter.
\end{grammarlens}

\subsection*{Your turn}
\begin{itemize}
\raggedright
  \item \emph{Say it:} \emph{tavo}
  \item \emph{Say it:} \emph{tavo}, once more
  \item \emph{Say it:} say \emph{tavo}
  \item \emph{Recall:} say the Marwadi for because, then the Marwadi for a hearth, then say \emph{tavo} again
\end{itemize}

\begin{tcolorbox}[breakable,colback=teal!4,colframe=teal!35!black,title={Before you move on}]
What does \mw{तवो} mean? (\textquotedblleft{}A griddle\textquotedblright{}.) Say it once more.
\end{tcolorbox}
\section[\texorpdfstring{\mw{कड़ाही} (kaṛāhī)}{कड़ाही (kaṛāhī)}]{\mw{कड़ाही} (kaṛāhī) — a wok}

\label{lesson:MW-C118-karahi}



\begin{quote}
Before the new one: say the Marwadi for a hearth, then the Marwadi for a griddle.
\end{quote}

\subsection*{You'll want to know: \mw{कड़ाही}}
\textbf{\mw{कड़ाही}} — \emph{kaṛāhī} — \textquotedblleft{}a wok\textquotedblright{}.

\begin{cousinweb}
\textbf{\mw{कड़ाही}} (\emph{kaṛāhī}) is a tadbhava, a word inherited from Sanskrit and reshaped by sound change. It comes from \emph{kaṭāha}, \textquotedblleft{}boiler, saucepan with handles\textquotedblright{}. CDIAL gives no Marwari form in this entry (CDIAL 2638); the match is through Hindi \emph{kaṛāhī}.
\end{cousinweb}

\begin{grammarlens}[title={What you've built}]
One more word for this chapter.
\end{grammarlens}

\subsection*{Your turn}
\begin{itemize}
\raggedright
  \item \emph{Say it:} \emph{kaṛāhī}
  \item \emph{Say it:} \emph{kaṛāhī}, once more
  \item \emph{Say it:} say \emph{kaṛāhī}
  \item \emph{Recall:} say the Marwadi for a hearth, then the Marwadi for a griddle, then say \emph{kaṛāhī} again
\end{itemize}

\begin{tcolorbox}[breakable,colback=teal!4,colframe=teal!35!black,title={Before you move on}]
What does \mw{कड़ाही} mean? (\textquotedblleft{}A wok\textquotedblright{}.) Say it once more.
\end{tcolorbox}
\section[\texorpdfstring{\mw{हांडी} (hā̃ḍī)}{हांडी (hā̃ḍī)}]{\mw{हांडी} (hā̃ḍī) — a cooking pot}

\label{lesson:MW-C118-hadi}



\begin{quote}
Before the new one: say the Marwadi for a griddle, then the Marwadi for a wok.
\end{quote}

\subsection*{You'll want to know: \mw{हांडी}}
\textbf{\mw{हांडी}} — \emph{hā̃ḍī} — \textquotedblleft{}a cooking pot\textquotedblright{}.

\begin{cousinweb}
\textbf{\mw{हांडी}} (\emph{hā̃ḍī}) is a tadbhava, a word inherited and reshaped by sound change. It comes from \emph{*hāṇḍa}, \textquotedblleft{}pot\textquotedblright{}, a form CDIAL reconstructs from its descendants. CDIAL gives no Marwari form in this entry (CDIAL 14050); the match is through Hindi \emph{hā̃ḍī} and Gujarati \emph{hā̃ḍī}. CDIAL says any link with Sanskrit \emph{bhāṇḍa}, \textquotedblleft{}pot\textquotedblright{}, is not clear.
\end{cousinweb}

\begin{grammarlens}[title={What you've built}]
One more word for this chapter.
\end{grammarlens}

\subsection*{Your turn}
\begin{itemize}
\raggedright
  \item \emph{Say it:} \emph{hā̃ḍī}
  \item \emph{Say it:} \emph{hā̃ḍī}, once more
  \item \emph{Say it:} say \emph{hā̃ḍī}
  \item \emph{Recall:} say the Marwadi for a griddle, then the Marwadi for a wok, then say \emph{hā̃ḍī} again
\end{itemize}

\begin{tcolorbox}[breakable,colback=teal!4,colframe=teal!35!black,title={Before you move on}]
What does \mw{हांडी} mean? (\textquotedblleft{}A cooking pot\textquotedblright{}.) Say it once more.
\end{tcolorbox}
\section[\texorpdfstring{\mw{घड़ो} (ghaṛo)}{घड़ो (ghaṛo)}]{\mw{घड़ो} (ghaṛo) — a water pot}

\label{lesson:MW-C118-gharo}



\begin{quote}
Before the new one: say the Marwadi for a wok, then the Marwadi for a cooking pot.
\end{quote}

\subsection*{You'll want to know: \mw{घड़ो}}
\textbf{\mw{घड़ो}} — \emph{ghaṛo} — \textquotedblleft{}a water pot\textquotedblright{}.

\begin{cousinweb}
\textbf{\mw{घड़ो}} (\emph{ghaṛo}) is a tadbhava, a word inherited from Sanskrit and reshaped by sound change. It comes from \emph{ghaṭa}, \textquotedblleft{}pot\textquotedblright{}. CDIAL lists Marwari \emph{ghaṛī} under it (CDIAL 4406). The Marwari row is \emph{ghaṛī}, \textquotedblleft{}water-clock\textquotedblright{}; Hindi \emph{ghaṛā} and Gujarati \emph{ghaṛo} mean \textquotedblleft{}water pot\textquotedblright{}, as \textbf{\mw{घड़ो}} (\emph{ghaṛo}) does.
\end{cousinweb}

\begin{grammarlens}[title={What you've built}]
One more word for this chapter.
\end{grammarlens}

\subsection*{Your turn}
\begin{itemize}
\raggedright
  \item \emph{Say it:} \emph{ghaṛo}
  \item \emph{Say it:} \emph{ghaṛo}, once more
  \item \emph{Say it:} say \emph{ghaṛo}
  \item \emph{Recall:} say the Marwadi for a wok, then the Marwadi for a cooking pot, then say \emph{ghaṛo} again
\end{itemize}

\begin{tcolorbox}[breakable,colback=teal!4,colframe=teal!35!black,title={Before you move on}]
What does \mw{घड़ो} mean? (\textquotedblleft{}A water pot\textquotedblright{}.) Say it once more.
\end{tcolorbox}
